\documentclass[10pt,a4paper]{amsart}
\usepackage[margin=19mm]{geometry}
\usepackage{amsmath,amssymb,mathtools}
\usepackage[scr=rsfs,cal=euler]{mathalfa}
\usepackage{xfrac}
\usepackage[unicode,hidelinks,bookmarks=false]{hyperref}
\newcommand{\ie}{i.e.\ }
\newcommand{\cf}{cf.\ }
\newcommand{\resp}{respectively\ }
\newcommand{\Subsection}{\subsection}
\providecommand{\nobreakdash}{\nobreak\hbox{-}}
\setlength{\emergencystretch}{2em}
\multlinegap0pt
\usepackage{bm}
\usepackage{bbm}
\usepackage{dsfont}
\usepackage{xspace}
\usepackage{cancel}
\usepackage{mathtools}
\usepackage{amsthm}
\makeatletter
\@ifundefined{comment}{\newtheorem{comment}{Comment}}{}
\@ifundefined{theorem}{\newtheorem{theorem}{Theorem}}{}
\@ifundefined{lemma}{\newtheorem{lemma}[theorem]{Lemma}}{}
\@ifundefined{proposition}{\newtheorem{proposition}[theorem]{Proposition}}{}
\makeatother
\newcommand{\F}{\mathbb{F}}
\newcommand{\Z}{\mathbb{Z}}
\newcommand{\bfe}{{\boldsymbol e}}
\newcommand{\bfg}{{\boldsymbol g}}
\title[Serving Every Symbol: All-Symbol PIR and Batch Codes]{Serving Every Symbol: All-Symbol PIR and Batch Codes}
\author{Avital Boruchovsky, Anina Gruica, Jonathan Niemann, Eitan Yaakobi}
\date{Source: arXiv:2601.04041v2 (CC BY 4.0)}
\begin{document}
\maketitle

\noindent\emph{Source and licence.} Serving Every Symbol: All-Symbol PIR and Batch Codes. Version arXiv:2601.04041v2 (CC BY 4.0), distributed under the Creative Commons Attribution 4.0 International licence, as recorded in the arXiv licence metadata for this exact version and shown on the abstract page https://arxiv.org/abs/2601.04041v2. Full rights notice: licenses/arxiv-2601.04041-CC-BY-4.0.txt.\par\smallskip

\noindent\textit{Editorial scope.} Counted unit: the Proposition whose source label is \texttt{prop:t=4\_q=2} in the article (source0.tex lines 895--901, source label \texttt{prop:t=4\_q=2}), delivered together with the article's own complete original proof of it (source0.tex lines 903--963, one \texttt{proof} environment of 61 lines). No mathematics is written, repaired or re-derived: statement and proof are the article's own text line for line. Required context retained uncounted: Lem:G'IsBatch (lines 763--765); thm:t=4 (lines 881--886). Every definition, display and label of those lines is retained unchanged. Omitted from this package and not redistributed: 1--762, 766--880, 887--894, 902--902, 964--1218. Nothing inside a retained excerpt is omitted or altered: every retained body is delivered exactly as the article's lines are written, so no omission receipt of any kind accompanies this package. Citations: the retained excerpts carry no citation command at all, so no bibliography entry of the article is transcribed and no reference is redistributed. Reference adaptation: none was needed and none was made; every reference inside the delivered unit resolves inside it, so no unresolved reference or citation is produced. Printed slips kept literal: no printed word, symbol, hypothesis or number of the counted statement or of its proof was altered. Layout and class: the article's own class is \texttt{article} with packages including amssymb, amsmath, xcolor, float, latexsym, amsthm, empheq, bm, booktabs, subcaption, tikz, pgfplots, cite, balance; the wrapper is the lane's pinned \texttt{amsart} editorial wrapper at 10pt on A4, which restates the theorem-like environments the retained text needs and adds the article's own macro definitions that the retained text needs (\texttt{\textbackslash F}, \texttt{\textbackslash Z}, \texttt{\textbackslash bfe}, \texttt{\textbackslash bfg}), with extra wrapper packages (bm, bbm, dsfont, xspace, cancel, mathtools). The delivered numbering is the unit's own. Assets: no retained span contains a figure environment, a graphics-inclusion command, a PDF-page inclusion, a table or a TikZ picture; no image file is copied into this package, the package declares no asset dependency and no image file is redistributed with it. Licence: version arXiv:2601.04041v2 of this article is distributed under the Creative Commons Attribution 4.0 International licence, as recorded in the arXiv licence metadata for this exact version and shown on the abstract page https://arxiv.org/abs/2601.04041v2. A full rights notice is delivered at licenses/arxiv-2601.04041-CC-BY-4.0.txt. Characters: every retained excerpt is pure ASCII, so the delivered engine, XeLaTeX with its default Latin Modern text font, reports no missing character.
\begin{lemma}\label{Lem:G'IsBatch}
    The matrix $G'$ can serve any list of four columns from $I_k$.
\end{lemma}

\begin{theorem} \label{thm:t=4} 
    For odd $q$, we have
    $$
        k+1 + \left\lceil \frac{1 + \sqrt{1 + 8k}}{2}\right\rceil \leq ASP(k,4,q) \leq ASB(k,4,q) \leq k+2 + \left\lceil \frac{1 + \sqrt{1 + 8k}}{2}\right\rceil.
    $$
\end{theorem}

\begin{proposition} \label{prop:t=4_q=2} 
    For $k \in S =\{1,2,3,4,5,7,8,11,12,16\}$ and $\ell \in \Z_{>0}$ we have 
    $$
        ASP(k,4,2^\ell)=ASB(k,4,2^\ell)= k+1 + \left\lceil \frac{1 + \sqrt{1 + 8k}}{2}\right\rceil,
    $$
    i.e., the lower bound from Theorem \ref{thm:t=4} is the true value for $k \in S$ when $q$ is even. 
\end{proposition}

\begin{proof}
    First, we prove that the parity column of $G'$ has four disjoint recovery sets if and only if $k \in S$. As we can take the column itself as one of the recovery sets, we need to find three additional recovery sets among the columns of $G$. Since each row of $G$ has weight exactly three, every column must belong to some recovery set and any two columns with overlapping support must be in disjoint recovery sets. 
    
    Consequently, the columns of $A$ must be partitioned into three sets $A_1, A_2, A_3$ such that no two columns within the same set overlap in support. If such a partition exists, then supplementing each $A_i$ with appropriate columns of $I_k$ yields three disjoint recovery sets for the parity column.
    
    Let $|A_i| = r_i$. Then
    $$
        r_1 + r_2 + r_3 = r = \left\lceil \frac{1 + \sqrt{1 + 8k}}{2}\right\rceil.
    $$
    
    For any of the $\binom{r_i}{2}$ pairs of columns in $A_i$ there is a weight-$2$ vector from $\F_2^r$ that cannot be used as a row in the construction of $A$ (as otherwise, the support of these two rows will overlap). This means that we must have 
    $$
        \binom{r}{2} - \left( \binom{r_1}{2} + \binom{r_2}{2} + \binom{r_3}{2}\right) \geq k,
    $$

    which can be rewritten as 
    $$
       r^2 - (r_1^2 + r_2^2 + r_3^2) \geq 2k.
    $$

    By the QM-AM inequality this implies 
    $$
        \frac{1}{3} \left\lceil \frac{1 + \sqrt{1 + 8k}}{2}\right\rceil^2 = \frac{1}{3}r^2 \geq k,
    $$
    and one can check that this holds only for $k \in S$. 

    Finally, for each $k \in S$ it can easily be checked that one can choose suitable subsets of columns of $A$.

    Next, we show that there are also four disjoint recovery sets for each column from $A$. In fact, we will show that there is a parition of the columns of $G'$ into four disjoint sets such that each of these is a recovery set, like for the columns of $I_k$. 

    
    Suppose $\bfg$ is a column of $A$.
    By the construction of $A$, the sum of all the other columns of $A$ is equal to $g$ (here we use the fact that $q$ has characteristic $2$).  Hence, one may take as recovery sets: the column $\bfg$ itself; all other columns of $A$; the columns of $I_k$ whose supports intersect the support of $\bfg$; and, finally, the remaining columns of  $=I_k$ together with the parity column.   
    \begin{comment}
        an argument similar to the one used in the proof of Lemma \ref{Lem:G'IsBatch} we may assume 
    $$
        \bfg = \bfg_{i_0} = \bfg_{i_1} + \sum_{i \in R_1} \bfe_i = \sum_{j \in R_2} \bfe_j,
    $$
    with $i_0,i_1 \in [n]\setminus [k]$ and $R_1,R_2 \subseteq [k]$. Again, let $T$ be the remaining indices of columns of~$G$. Then, 
    \begin{align*}
        0 & =  \bfg_{n+1} + \sum_{i\in R_1} \bfg_i + \sum_{j\in R_2}\bfg_j + \sum_{l\in R_3} \bfg_l + \sum_{m\in T}\bfg_m \\
        &= \bfg_{n+1} + \bfg + \bfg + \bfg  + \sum_{m\in T}\bfg_m \\
        &= \bfg_{n+1} + \bfg + \sum_{m\in T}\bfg_m,
    \end{align*}
    so that 
    $$
    \bfg  = \bfg_{n+1}+ \sum_{m\in T}\bfg_m,
    $$
    and the claim follows. \newline
\end{comment}

    To go from the all-symbol PIR property to the all-symbol batch property, the only non-trivial thing we need to check is that $G'$ can serve a list with two copies of two distinct columns, not both in $I_k$. If the parity column does not appear, then we can use the same argument as in the proof of Lemma \ref{Lem:G'IsBatch}. So, suppose instead that there are two copies of the parity check column together with two copies of another column, say $\bfg_{i_1}$. 
    
    We know that $\bfg_{i_1}$ has four recovery sets that form a partition of the columns of $G'$. Let $R_1,R_2,R_3,R_4\subseteq [n+1]$ be the sets of indices that correspond to these recovery sets. We may assume, without loss of generality, that $R_1 = \{i_1\}$ and that $n+1\in R_4$.

    Now, 
    $$
        \bfg_{n+1} = \sum_{i \in [n]} \bfg_i  = \bfg_{i_0} + \sum_{i\in R_2}\bfg_i + \sum_{j\in R_3}\bfg_j + \sum_{l \in R_4\setminus\{n+1\}} \bfg_l = 3\cdot\sum_{j\in R_3}\bfg_j + \sum_{l \in R_4\setminus\{n+1\}} \bfg_l,
    $$
    so we can take the columns corresponding to $R_1$ and $R_2$ as recovery sets for $\bfg_{i_1}$, and $\{n+1\}$ and $R_3 \cup R_4\setminus \{n+1\}$ as recovery sets for $\bfg_{n+1}$. This finishes the proof.
\end{proof}

\end{document}
